% GENERATED FILE. Edit canonical lessons, then run npm run generate:books.
% canonical-source-hash: fnv1a64:10d036b8783c1c4c
% canonical-lessons: SA-C389-sakyah, SA-C389-sadhyah, SA-C389-kalpayati, SA-C389-kalpana, SA-C389-yatharthah

\chapter{Can It Be Done?}
\label{ch:sa-can-it-be-done}
\hlchaptermodality{389}

\begin{chapteropening}
\textbf{By the end of this chapter:} \emph{I can say whether a plan is possible and what we imagine might happen.}

Complete the last lesson of chapter 389: I can say whether a plan is possible and what we imagine might happen.
\end{chapteropening}

\section[\texorpdfstring{\sk{शक्यः} (śakyaḥ)}{शक्यः (śakyaḥ)}]{\sk{शक्यः} (śakyaḥ) — possible; able to be done}

\label{lesson:SA-C389-sakyah}



\begin{quote}
Before the new one: say the Sanskrit for in vain, then the Sanskrit for an ideal.
\end{quote}

\subsection*{You'll want to know: \sk{शक्यः}}
\textbf{\sk{शक्यः}} — \emph{śakyaḥ} — \textquotedblleft{}possible; able to be done\textquotedblright{}.

\begin{quote}
\sk{यदि} \sk{समयः} \sk{अस्ति}, \sk{तर्हि} \sk{एतत्} \sk{कार्यं} \sk{शक्यम्।}
\end{quote}

\begin{quote}
\emph{yadi samayaḥ asti, tarhi etat kāryaṃ śakyam.}
\end{quote}

\begin{quote}
If there is time, this work can be done.
\end{quote}

\begin{cousinweb}
\textbf{\sk{शक्यः}} (\emph{śakyaḥ}) comes from the root √śak, \textquotedblleft{}to be able\textquotedblright{}; MW gives \emph{śakya} \textquotedblleft{}able, possible, practicable\textquotedblright{} (MW).
\end{cousinweb}

\begin{grammarlens}[title={What you've built}]
One more word for saying what happens if.
\end{grammarlens}

\subsection*{Your turn}
\begin{itemize}
\raggedright
  \item \emph{Say it:} \emph{śakyaḥ}
  \item \emph{Say it:} \emph{śakyaḥ}, once more
  \item \emph{Say it:} \emph{yadi samayaḥ asti, tarhi etat kāryaṃ śakyam.} — If there is time, this work can be done.
  \item \emph{Recall:} say the Sanskrit for in vain, then the Sanskrit for an ideal, then say \emph{śakyaḥ} again
\end{itemize}

\begin{tcolorbox}[breakable,colback=teal!4,colframe=teal!35!black,title={Before you move on}]
What does \sk{शक्यः} mean? (\textquotedblleft{}Possible; able to be done\textquotedblright{}.) Say it once more.
\end{tcolorbox}
\section[\texorpdfstring{\sk{साध्यः} (sādhyaḥ)}{साध्यः (sādhyaḥ)}]{\sk{साध्यः} (sādhyaḥ) — achievable}

\label{lesson:SA-C389-sadhyah}



\begin{quote}
Before the new one: say the Sanskrit for an ideal, then the Sanskrit for possible.
\end{quote}

\subsection*{You'll want to know: \sk{साध्यः}}
\textbf{\sk{साध्यः}} — \emph{sādhyaḥ} — \textquotedblleft{}achievable\textquotedblright{}.

\begin{quote}
\sk{यदि} \sk{योजना} \sk{उत्तमा}, \sk{तर्हि} \sk{लक्ष्यं} \sk{साध्यम्।}
\end{quote}

\begin{quote}
\emph{yadi yojanā uttamā, tarhi lakṣyaṃ sādhyam.}
\end{quote}

\begin{quote}
If the plan is good, the goal is achievable.
\end{quote}

\begin{cousinweb}
\textbf{\sk{साध्यः}} (\emph{sādhyaḥ}) comes from the root √sidh, \textquotedblleft{}to succeed\textquotedblright{}, in the stem \emph{sādh-}; MW gives \emph{sādhya} \textquotedblleft{}to be accomplished\textquotedblright{} and \textquotedblleft{}practicable, feasible, attainable\textquotedblright{} (MW).
\end{cousinweb}

\begin{grammarlens}[title={What you've built}]
One more word for saying what happens if.
\end{grammarlens}

\subsection*{Your turn}
\begin{itemize}
\raggedright
  \item \emph{Say it:} \emph{sādhyaḥ}
  \item \emph{Say it:} \emph{sādhyaḥ}, once more
  \item \emph{Say it:} \emph{yadi yojanā uttamā, tarhi lakṣyaṃ sādhyam.} — If the plan is good, the goal is achievable.
  \item \emph{Recall:} say the Sanskrit for an ideal, then the Sanskrit for possible, then say \emph{sādhyaḥ} again
\end{itemize}

\begin{tcolorbox}[breakable,colback=teal!4,colframe=teal!35!black,title={Before you move on}]
What does \sk{साध्यः} mean? (\textquotedblleft{}Achievable\textquotedblright{}.) Say it once more.
\end{tcolorbox}
\section[\texorpdfstring{\sk{कल्पयति} (kalpayati)}{कल्पयति (kalpayati)}]{\sk{कल्पयति} (kalpayati) — imagines; supposes}

\label{lesson:SA-C389-kalpayati}



\begin{quote}
Before the new one: say the Sanskrit for possible, then the Sanskrit for achievable.
\end{quote}

\subsection*{You'll want to know: \sk{कल्पयति}}
\textbf{\sk{कल्पयति}} — \emph{kalpayati} — \textquotedblleft{}imagines; supposes\textquotedblright{}.

\begin{quote}
\sk{यदि} \sk{बालः} \sk{एकाकी}, \sk{तर्हि} \sk{सः} \sk{मित्राणि} \sk{कल्पयति।}
\end{quote}

\begin{quote}
\emph{yadi bālaḥ ekākī, tarhi saḥ mitrāṇi kalpayati.}
\end{quote}

\begin{quote}
If the child is alone, he imagines friends.
\end{quote}

Past (he …): \sk{सः} \sk{नूतनं} \sk{जगत्} \sk{कल्पितवान्}, \emph{saḥ nūtanaṃ jagat kalpitavān}.

\begin{cousinweb}
\textbf{\sk{कल्पयति}} (\emph{kalpayati}) comes from the root √kḷp, \textquotedblleft{}to be fit\textquotedblright{}, in its causative form; MW gives the senses \textquotedblleft{}to frame, to form, to invent\textquotedblright{} and \textquotedblleft{}to imagine\textquotedblright{} (MW).
\end{cousinweb}

\begin{grammarlens}[title={What you've built}]
One more word for saying what happens if.
\end{grammarlens}

\subsection*{Your turn}
\begin{itemize}
\raggedright
  \item \emph{Say it:} \emph{kalpayati}
  \item \emph{Say it:} \emph{kalpayati}, once more
  \item \emph{Say it:} \emph{yadi bālaḥ ekākī, tarhi saḥ mitrāṇi kalpayati.} — If the child is alone, he imagines friends.
  \item \emph{Recall:} say the Sanskrit for possible, then the Sanskrit for achievable, then say \emph{kalpayati} again
\end{itemize}

\begin{tcolorbox}[breakable,colback=teal!4,colframe=teal!35!black,title={Before you move on}]
What does \sk{कल्पयति} mean? (\textquotedblleft{}Imagines; supposes\textquotedblright{}.) Say it once more.
\end{tcolorbox}
\section[\texorpdfstring{\sk{कल्पना} (kalpanā)}{कल्पना (kalpanā)}]{\sk{कल्पना} (kalpanā) — imagination; a fancy}

\label{lesson:SA-C389-kalpana}



\begin{quote}
Before the new one: say the Sanskrit for achievable, then the Sanskrit for imagines.
\end{quote}

\subsection*{You'll want to know: \sk{कल्पना}}
\textbf{\sk{कल्पना}} — \emph{kalpanā} — \textquotedblleft{}imagination; a fancy\textquotedblright{}.

\begin{quote}
\sk{यदि} \sk{कल्पना} \sk{न} \sk{अभविष्यत्}, \sk{तर्हि} \sk{कविः} \sk{काव्यं} \sk{न} \sk{अकरिष्यत्।}
\end{quote}

\begin{quote}
\emph{yadi kalpanā na abhaviṣyat, tarhi kaviḥ kāvyaṃ na akariṣyat.}
\end{quote}

\begin{quote}
Had there been no imagination, the poet would have made no poem.
\end{quote}

A feminine noun.

\begin{cousinweb}
\textbf{\sk{कल्पना}} (\emph{kalpanā}) comes from the root √kḷp, \textquotedblleft{}to be fit, to form\textquotedblright{}; MW gives \emph{kalpanā} \textquotedblleft{}creating in the mind, assuming anything to be real, fiction\textquotedblright{} (MW).
\end{cousinweb}

\begin{grammarlens}[title={What you've built}]
One more word for saying what happens if.
\end{grammarlens}

\subsection*{Your turn}
\begin{itemize}
\raggedright
  \item \emph{Say it:} \emph{kalpanā}
  \item \emph{Say it:} \emph{kalpanā}, once more
  \item \emph{Say it:} \emph{yadi kalpanā na abhaviṣyat, tarhi kaviḥ kāvyaṃ na akariṣyat.} — Had there been no imagination, the poet would have made no poem.
  \item \emph{Recall:} say the Sanskrit for achievable, then the Sanskrit for imagines, then say \emph{kalpanā} again
\end{itemize}

\begin{tcolorbox}[breakable,colback=teal!4,colframe=teal!35!black,title={Before you move on}]
What does \sk{कल्पना} mean? (\textquotedblleft{}Imagination; a fancy\textquotedblright{}.) Say it once more.
\end{tcolorbox}
\section[\texorpdfstring{\sk{यथार्थः} (yathārthaḥ)}{यथार्थः (yathārthaḥ)}]{\sk{यथार्थः} (yathārthaḥ) — true to fact; real}

\label{lesson:SA-C389-yatharthah}



\begin{quote}
Before the new one: say the Sanskrit for imagines, then the Sanskrit for imagination.
\end{quote}

\subsection*{You'll want to know: \sk{यथार्थः}}
\textbf{\sk{यथार्थः}} — \emph{yathārthaḥ} — \textquotedblleft{}true to fact; real\textquotedblright{}.

\begin{quote}
\sk{यदि} \sk{वार्ता} \sk{यथार्था}, \sk{तर्हि} \sk{सः} \sk{श्वः} \sk{गमिष्यति।}
\end{quote}

\begin{quote}
\emph{yadi vārtā yathārthā, tarhi saḥ śvaḥ gamiṣyati.}
\end{quote}

\begin{quote}
If the news is true, he will go tomorrow.
\end{quote}

\begin{cousinweb}
\textbf{\sk{यथार्थः}} (\emph{yathārthaḥ}) is built from \emph{yathā}, \textquotedblleft{}as\textquotedblright{}, and \emph{artha}, \textquotedblleft{}meaning\textquotedblright{}; MW gives \emph{yathārtha} \textquotedblleft{}accordant with reality, conformable to truth, true, genuine\textquotedblright{} (MW).
\end{cousinweb}

\begin{grammarlens}[title={What you've built}]
One more word for saying what happens if.
\end{grammarlens}

\subsection*{Your turn}
\begin{itemize}
\raggedright
  \item \emph{Say it:} \emph{yathārthaḥ}
  \item \emph{Say it:} \emph{yathārthaḥ}, once more
  \item \emph{Say it:} \emph{yadi vārtā yathārthā, tarhi saḥ śvaḥ gamiṣyati.} — If the news is true, he will go tomorrow.
  \item \emph{Recall:} say the Sanskrit for imagines, then the Sanskrit for imagination, then say \emph{yathārthaḥ} again
\end{itemize}

\begin{tcolorbox}[breakable,colback=teal!4,colframe=teal!35!black,title={Before you move on}]
What does \sk{यथार्थः} mean? (\textquotedblleft{}True to fact; real\textquotedblright{}.) Say it once more.
\end{tcolorbox}
